\subsection{微分方程解空间上的高阶对称}

高自旋场的规范约束以前，先区分使场方程解空间保持不变的微分对称与普通算符对易。

一般流形上的高自旋守恒流之所以困难，并不是缺乏排列指标的技巧，而是微分算子的对称性在曲率下受到阻碍。最干净的数学入口是：哪些高阶微分算子能将自由场方程的解送到另一个解？这可以把高自旋守恒问题转化为主符号、共形 Killing 张量和低阶曲率补偿项之间的相容条件。

设 $L$ 是一个二阶 Laplace 型算子，$D$ 是一个阶数为 $S$ 的微分算子。局部写为
\[
 D=K^{a_1\cdots a_S}(x)\nabla_{a_1}\cdots\nabla_{a_S}
 +\text{低阶导数项}.
\]
其最高阶系数的对称部分可看作 $T^*M$ 上次数为 $S$ 的多项式
\[
 \sigma_S(D)(x,p)=K^{a_1\cdots a_S}(x)p_{a_1}\cdots p_{a_S}.
\]
主符号独立于所取局部坐标；因此它比微分算子的具体展开更适合作为几何研究对象。

\begin{definition}[模场方程的高阶对称算子]
给定二阶算子 $L$，若存在微分算子 $D'$ 使
\begin{equation}\label{eq:ghs-on-shell-symmetry}
 L D=D'L,
\end{equation}
则称 $D$ 为 $L$ 的一个解空间对称算子。于是 $Lf=0$ 必推出 $L(Df)=0$。
\end{definition}

条件 \cref{eq:ghs-on-shell-symmetry} 比 $[L,D]=0$ 更一般：后者要求在所有函数上对易，前者只要求 $D$ 维持解空间。平直共形场中的很多高阶对称性正属于这一类。

考虑 $L=-\Delta_g+\xi_d\mathcal R_g$ 的主符号 $\ell(x,p)=g^{ab}p_ap_b$。把 $LD-D'L$ 的最高一阶非平凡导数项提取出来，等价于分析主符号的 Poisson 括号 $\{\ell,\sigma_S(D)\}$ 模主符号 $\ell$ 所生成的多项式理想的等价类。其迹自由部分为
\begin{equation}\label{eq:ghs-ck-tensor-principal}
 \operatorname{STF}\nabla_{(a_0}K_{a_1\cdots a_S)}=0.
\end{equation}
这里比较的是主符号的多项式可除性：$\{\ell,\sigma_S(D)\}$ 必须是 $\ell$ 的倍数。正定度规的实零锥只有零点，不能以实零锥上的点值条件代替这一多项式条件。取最高次系数的无迹部分就得到上式。
这是一种必要的共形 Killing 张量条件；对于一般曲率背景，它不必是存在完整对称算子的充分条件。

自旋一情形的主符号可以直接在局部坐标中检验。
令 $D=\xi^a\nabla_a+\alpha\nabla_a\xi^a$。暂不考虑低阶曲率项，直接计算
\begin{align*}
[-\Delta,\xi^a\nabla_a]f
&=-(\Delta\xi^a)\nabla_af
 -2(\nabla^b\xi^a)\nabla_b\nabla_af
 -\xi^a[\Delta,\nabla_a]f.
\end{align*}
二阶 Hessian 对称，因此最高导数项只看 $\nabla^{(b}\xi^{a)}$。若仅要求 $D$ 将 $Lf=0$ 的解映回解，允许这项与 $Lf$ 的主符号成比例，即
$\nabla_{(a}\xi_{b)}=\frac1d(\nabla\cdot\xi)g_{ab}$。
这正是 $S=1$ 的 \cref{eq:ghs-ck-tensor-principal}。低阶项必须再通过 $\alpha$ 与曲率修正配平；不能仅靠最高阶方程证明对称算子存在。

在平直空间，完整共形代数包含平移、转动、尺度与特殊共形变换。这些一阶对称作用可以相乘形成高阶微分算子，并在适当的低阶修正与解空间等价关系下组织为非平凡的高阶对称代数。Eastwood 对欧氏 拉普拉斯算子 的高阶对称进行了系统分类，并讨论了一般共形流形上的相应几何结构 \cite{eastwood-higher}。

当曲率不为零且随位置变化时，两个协变导数的交换会不断生成 $R\cdot\nabla^{S-2}$ 及其导数。要保持 \cref{eq:ghs-on-shell-symmetry}，必须逐阶补充低阶系数并检验可解性。即使找到了形式上的共形 Killing 张量，也可能被后续的曲率兼容条件阻挡。因此高自旋流守恒不是“给一般度规添加联络”就能自动修复的。

令 $\mathfrak g$ 为几何上实际存在的 Killing 或共形 Killing 代数。一阶生成元作用在自由场解空间上，形成一个 $\mathfrak g$ 表示 $\rho$。其泛包络代数 $U(\mathfrak g)$ 通过乘积得到更高阶算子，但不同包络元素可能在场方程的解空间上产生相同作用；因此真正的对称代数通常是 $U(\mathfrak g)$ 对某个双边理想的商。

这与群论中“表示不一定忠实”是同一个数学现象：把一组抽象代数元素作用到一个具体空间时，必须将作用为零的部分作为核商掉。平直共形 拉普拉斯算子的高阶对称代数提供了这类商代数的非平凡实例。一般背景的 $\mathfrak g$ 可能很小甚至为零，但局部微分算子和其谱仍然有意义。

高阶对称算子的存在可以使自由场的相关函数满足额外 Ward 恒等式；若将其对应的二次型高自旋流构造成局域复合算符，还需要给出正规序、守恒改进项及量子异常检查。相互作用改变运动方程，因而往往打破这些额外高阶对称性。后文平直 共形场论 的高自旋反常维数计算，就是对这一几何事实的一个动力学测量：经典非守恒项进入量子两点函数，守恒短表示一般重组为长表示。
